\documentclass[10pt,a4paper]{amsart}
\usepackage[margin=19mm]{geometry}
\usepackage{amsmath,amssymb,mathtools}
\usepackage[scr=rsfs,cal=euler]{mathalfa}
\usepackage{xfrac}
\usepackage[unicode,hidelinks,bookmarks=false]{hyperref}
\newcommand{\ie}{i.e.\ }
\newcommand{\cf}{cf.\ }
\newcommand{\resp}{respectively\ }
\newcommand{\Subsection}{\subsection}
\providecommand{\nobreakdash}{\nobreak\hbox{-}}
\setlength{\emergencystretch}{2em}
\multlinegap0pt
\usepackage{dsfont}
\usepackage{float}
\usepackage{xcolor}
\usepackage{amssymb}
\usepackage{enumerate}
\usepackage{tikz}
\usepackage{csquotes}
\usepackage{mathtools}
\newtheorem{lemma}{Lemma}
\DeclareMathOperator*{\Res}{Res}
\title{Les Houches lecture notes on topological recursion}
\author{Vincent Bouchard}
\date{Source: arXiv:2409.06657v1 (CC BY 4.0)}
\begin{document}
\maketitle

\noindent\emph{Source and licence.} Les Houches lecture notes on topological recursion. Version arXiv:2409.06657v1 (CC BY 4.0), distributed by arXiv under the Creative Commons Attribution 4.0 International licence, http://creativecommons.org/licenses/by/4.0/, as recorded in the arXiv licence metadata for this exact version and shown on the abstract page https://arxiv.org/abs/2409.06657v1. Full licence notice: \texttt{licenses/arxiv-2409.06657-CC-BY-4.0.txt}.\par\smallskip

\noindent\textit{Editorial scope.} Counted unit: the article's lemma lemma (source0.tex lines 848--850). The article's complete original proof of it is delivered with it (source0.tex lines 852--870, one \texttt{proof} environment of 19 lines). No mathematics is written, repaired or re-derived: the statement and the proof are the article's own lines, in the article's own notation, and no retained line is substituted or re-flowed.  No further source-local context is needed: every cross-reference of every retained line resolves inside the two delivered spans.  Nothing was omitted from any retained span, so the package carries no omission record.  No retained line contains a citation, so the wrapper carries no bibliography and no bibliography entry.  No retained line contains a figure environment, an image-inclusion command or drawing code, so no image is delivered and the package declares no asset dependency.  Editorial formatting only, in addition: an AMS \texttt{amsart} preamble that restates the article's own declarations and macros used by the retained lines, namely the environments \texttt{lemma} on separate counters, together with the standard packages those lines need.  The retained environments print in this document as Lemma 1 in order of appearance, whereas the article prints the same items with the numbering of its own full document.  The complete native source of the article as distributed in the arXiv e-print for this exact version is bundled unchanged as \texttt{sources/165063/source0.tex}.
\begin{lemma}
The projection $\hat{B}_P[\alpha]$ is a one-form on $\Sigma$ that has the same principal part as $\alpha$ on $P$. It then follows that it is indeed a projection, as $\hat{B}_P \circ \hat{B}_P = \hat{B}_P$.
\end{lemma}

\begin{proof}
This is easiest to see in local coordinates. Pick a point $p \in P$ and a local coordinate $\zeta$ centered at $p$. We can expand the fundamental bidifferential and $\alpha$ in this local coordinate:
\begin{align}
\omega_{0,2}\simeq& \frac{d \zeta_1 d \zeta_2}{(\zeta_1-\zeta_2)^2} + \sum_{l,m> 0} \phi_{lm}  \zeta_1^{l-1} \zeta_2^{m-1} d \zeta_1 d \zeta_2,\\
\alpha \simeq& \sum_{k=-M}^\infty \alpha_k \zeta^{k} d \zeta.
\end{align}
Evaluating the projection at  $p$ in the local coordinate, we get:
\begin{align}
\hat{B}_p[\alpha](\zeta') \simeq& \Res_{\zeta=0} \left( \int^\zeta_0 \omega_{0,2}(\zeta,\zeta') \right) \alpha(\zeta) \\
=& \left( \sum_{k=-M}^{-1} \alpha_k (\zeta')^k  + O(1) \right) d \zeta'.
\end{align}
In other words, the one-forms $\hat{B}_p[\alpha]$ and $\alpha$ have the same principal part at $p$. The same holds true for the one-form $\hat{B}_P[\alpha]$, where we now sum over all $q \in P$, since the projection at the other points $q \neq p$ in $P$ are holomorphic at $p \in P$.

As for showing that it is a projection, for any one-form $\alpha$ on $\Sigma$, we can write
\begin{equation}
\alpha = \hat{B}_P[\alpha] + \omega
\end{equation}
where $\omega$ is holomorphic on $P$. (This is clearly true since $\alpha$ and $\hat{B}_P[\alpha]$ have the same principal part on $P$.) But then, it is easy to see that $\hat{B}_P[\omega] = 0$, and hence $\hat{B}_P \circ \hat{B}_P = \hat{B}_P$.
\end{proof}

\end{document}
